\documentclass[10pt,a4paper]{amsart}
\usepackage[margin=19mm]{geometry}
\usepackage{amsmath,amssymb,mathtools}
\usepackage[scr=rsfs,cal=euler]{mathalfa}
\usepackage{xfrac}
\usepackage[unicode,hidelinks,bookmarks=false]{hyperref}
\newcommand{\ie}{i.e.\ }
\newcommand{\cf}{cf.\ }
\newcommand{\resp}{respectively\ }
\newcommand{\Subsection}{\subsection}
\providecommand{\nobreakdash}{\nobreak\hbox{-}}
\setlength{\emergencystretch}{2em}
\multlinegap0pt
\usepackage{amssymb}
\usepackage{float}
\usepackage{enumerate}
\usepackage{tikz}
\usepackage[shortlabels]{enumitem}
\usepackage{mathtools}
\newtheorem{lem}{Lemma}
\newtheorem{rem}{Remark}
\newtheorem{prop}{Proposition}
\newtheorem{thm}{Theorem}
\def\dd{{\rm d}}
\title{Supercritical mean--field limit for magnetized Hamiltonian dynamics}
\author{Immanuel Ben-Porat, Gui-Qiang G. Chen, Difan Yuan}
\date{Source: arXiv:2609.00598v1 (CC BY 4.0)}
\begin{document}
\maketitle

\noindent\emph{Source and licence.} Supercritical mean--field limit for magnetized Hamiltonian dynamics. Version arXiv:2609.00598v1 (CC BY 4.0), distributed by arXiv under the Creative Commons Attribution 4.0 International licence, http://creativecommons.org/licenses/by/4.0/, as recorded in the arXiv licence metadata for this exact version and shown on the abstract page https://arxiv.org/abs/2609.00598v1. Full licence notice: \texttt{licenses/arxiv-2609.00598-CC-BY-4.0.txt}.\par\smallskip

\noindent\textit{Editorial scope.} Counted unit: the article's thm thm (source0.tex lines 2121--2127). The article's complete original proof of it is delivered with it (source0.tex lines 2128--2244, one \texttt{proof} environment of 117 lines). No mathematics is written, repaired or re-derived: the statement and the proof are the article's own lines, in the article's own notation, and no retained line is substituted or re-flowed.  Retained uncounted as the source-local context the proof needs: lines 131--150; lines 181--195; lines 1552--1554; lines 1615--1622; lines 1642--1644; lines 1936--1962; lines 1965--1977; lines 1979--2043; lines 2045--2062; lines 2058--2060; lines 2250--2276.  Nothing was omitted from any retained span, so the package carries no omission record.  No retained line contains a citation, so the wrapper carries no bibliography and no bibliography entry.  No retained line contains a figure environment, an image-inclusion command or drawing code, so no image is delivered and the package declares no asset dependency.  Editorial formatting only, in addition: an AMS \texttt{amsart} preamble that restates the article's own declarations and macros used by the retained lines, namely the environments \texttt{lem}, \texttt{rem}, \texttt{prop}, \texttt{thm} on separate counters, together with the standard packages those lines need.  The retained environments print in this document as Lemma 1, Remark 1, Proposition 1, Theorem 1 in order of appearance, whereas the article prints the same items with the numbering of its own full document.  The complete native source of the article as distributed in the arXiv e-print for this exact version is bundled unchanged as \texttt{sources/209042/source0.tex}.
\subsection{The quasi-neutral and supercritical mean field limits}
Consider now the magnetized Vlasov-Poisson equation with the following scaling, known as the \textit{quasi-neutral scaling}:
\begin{align}\tag{VPB$\varepsilon$}
\begin{cases}
\partial_{t}f_{\varepsilon}+\xi\cdot \nabla_{x}f_{\varepsilon}-(\nabla_{x}\Phi_{\varepsilon}+ B\xi^{\perp})\cdot\nabla_{\xi} f_{\varepsilon}=0, \qquad\, f_{\varepsilon}(0,\cdot)=f^{\mathrm{in}}_{\varepsilon},\\
\varepsilon\Delta_{x} \Phi_{\varepsilon}=1-\rho_{\varepsilon},
\end{cases}
\label{magnetized vp with quasineutral scaling}
\end{align}
with $\rho_{\varepsilon}(t,x)=\int_{\mathbb{R}^{2}} f_{\varepsilon}(t,x,\xi)\ \dd \xi$. 
It is instructive at this stage to comment about the physical interpretation of the quantities involved in the above equation:
\begin{itemize}
     \item[{\rm(i)}] $\rho_{\varepsilon}$ is the density of  the electrons. 
   \item[{\rm(ii)}] $1$ in the Poisson equation
   is the stationary density of the ions. 
   \item[{\rm(iii)}] $\varepsilon$ is the electric permittivity constant, which measures how close the density of the electrons is to the density of the ions.  
\end{itemize}
Under uniform bounds that control $\varepsilon\Delta\Phi_\varepsilon$, one expects the electron density to approach the ion density as $\varepsilon\to0$. This limit regime is referred to as the quasi-neutral limit. 
Since plasmas with $0$ total charge exhibit a fluid-like behavior, 
the quasi-neutral limit can be thought of as a hydrodynamical limit. 

\smallskip
Our second main objective is to study the supercritical mean field limit, which can be viewed as a \textit{combined quasi-neutral mean field limit}. More specifically, consider the microscopic dynamics associated with \eqref{magnetized vp with quasineutral scaling}, that is, the following system of $2\times 2N$ ODEs: 
\vspace{2pt}
\begin{align}\tag{NB$\varepsilon$}
 \begin{cases}
\frac{\dd}{\dd t}x_{i}^{N}(t)=\xi_{i}^{N}(t), &\,\, x_{i}^{N}(0)=x_{i}^{N,0},\\[1mm]
\frac{\dd}{\dd t}\xi_{i}^{N}(t)=-B(t,x_{i}^{N}(t))(\xi_{i}^{N}(t))^{\perp}-\frac{1}{N\varepsilon}\underset{j:j\neq i}{\sum}\nabla_{x} V(x_{i}^{N}(t)-x_{j}^{N}(t)),
&\,\,\xi_{i}^{N}(0)=\xi_{i}^{N,0}. 
\end{cases} 
\label{trajectories intro}
\end{align}
Here $V$ is the Green kernel on the torus, {\it i.e.}, the solution of 
$$
\Delta_{x}V=1-\delta_{0} 
\quad \mbox{on $\mathbb{T}^{2}\qquad$  \text{with} $\int_{\mathbb{T}^{2}}V(x)\, \dd x=0$}.$$

\begin{align}
\mathcal{F}_{\varepsilon,N}(t):=\frac{1}{2N}\sum_{i=1}^{N}\left\vert \xi_{i}(t)\right\vert^{2}+\frac{1}{2N^{2}\varepsilon}\sum_{i\neq j}V_{\geq 0}(x_{i}(t)-x_{j}(t)),\label{total energy discrete}     
\end{align}

 \begin{lem} \label{Conservation of discrete energy}
Let $(x_{1}^{0},\cdots\!,x_{N}^{0})\in \Delta_{N}^{c}$, 
and let $(x_{1}(t),\cdots\!,x_{N}(t),\xi_{1}(t),\cdots\!,\xi_{N}(t))\in C^{1}(\mathbb{R};\mathbb{T}^{2N}\times \mathbb{R}^{2N})$ 
be a solution to \eqref{trajectories intro}. Let $\mathcal{F}_{\varepsilon,N}(t)$ be given by \eqref{total energy discrete}. Then
\begin{align*}
\frac{\dd}{\dd t}\mathcal{F}_{\varepsilon,N}(t)=0 \qquad \mbox{for all}\ t\in \mathbb{R}_{+}.     
\end{align*}
\end{lem}

\begin{rem}
The conservation of energy as stated in {\rm Lemma \ref{Conservation of discrete energy}} holds whenever $V$ is even, $C^1$, and has a bounded gradient. The utility of this remark will become evident in {\rm Section \ref{wellposed-trajectories-sec}}.   \label{rem about conse of energy}  
\end{rem}

\begin{align*}
\underset{t\in [0,T]}{\sup}\left\Vert (J_{\varepsilon,N}-v\,\dd x)(t,\cdot)\right\Vert_{\mathrm{BL}^{*}}\underset{N\rightarrow \infty}{\longrightarrow} 0,
\end{align*}
where the bounded Lipschitz dual norm is the one defined above.
\qed


\section{Well-Posedness of Problem \eqref{trajectories intro}}
\label{wellposed-trajectories-sec}
In this section, we establish the existence and uniqueness of global solutions to \eqref{trajectories intro}. The proof proceeds by first establishing uniform separation of the particles and then applying a blow-up alternative. The presence of a general external magnetic field requires additional pointwise control of the particle velocities, which ensures that the system remains globally Lipschitz along the flow (rather than only locally Lipschitz). Note also that the system is non-homogeneous since $B$ is time dependent. 
We begin by recalling that the $2$-D periodic Coulomb potential is a smooth local perturbation of its whole-space counterpart. More precisely,
$
V\in C^\infty(\mathbb T^2\setminus\{0\}),
$
and, defining
$
V_{\mathbb R^2}(x):=-\frac{1}{2\pi}\log|x|,
$
there exists an even function
$
V_{\mathrm{loc}}\in C^\infty\big((\mathbb R^2\setminus\mathbb Z^2)\cup\{0\}\big)
$
such that
\begin{align*}
V(x)
=V_{\mathbb{R}^{2}}(x)+V_{\mathrm{loc}}(x) \qquad \mbox{for all}\ x\in \mathbb{T}^{2}\setminus \{0\}.     
\end{align*}

\begin{prop}[Local well-posedness]\label{Short time well posed}
Suppose that
\begin{itemize}
\item[{\rm(i)}]$X_{N}^{0}=(x_{1}^{N,0},\cdots\!,x_{N}^{N,0})\in \Delta_{N}^{c}$ and $\Xi_{N}^{0}=(\xi_{1}^{N,0},\cdots\!,\xi_{N}^{N,0})\in \mathbb{R}^{2N};$
\smallskip
\item[{\rm(ii)}] $B\in C(\mathbb{R}_{+};W^{1,\infty}(\mathbb{T}^{2}))$. 
\end{itemize}
Then there exist $T_{\ast}>0$ such that \eqref{trajectories intro} 
has a unique solution 
$$
(X_{N}(t),\Xi_{N}(t))\in C^{1}([0,T^{\ast}];\mathbb{T}^{2N}\times \mathbb{R}^{2N}).
$$ 
\end{prop}

\begin{proof}
Let $\Psi\in C^{\infty}_{0}(\mathbb{R})$ be an even function such that $\Psi(r)=1$ for $\left| r \right|\leq \frac{1}{8}$ and $\Psi(r)=0$ for $\left| r \right|\geq \frac{1}{4}$. 
Let $\psi\in C^{\infty}(\mathbb{T}^{2})$ be the radial function defined 
by $\psi(x) := \Psi(\left| x \right|)$. Set  
\[
\psi_{\eta}(x):=\Psi(\frac{\left|x\right|^{2}}{\eta^{2}}), \qquad\ 
V_{\eta}(x):= V_{\mathrm{loc}}(x)+(1-\psi_{\eta}(x))V_{\mathbb{R}^{2}}(x). 
\]
For brevity, in what follows we omit superscript $N$ from $x_{i,\eta}^{N}$ and $\xi_{i,\eta}^{N}$. Consider the regularized system: 
\begin{align}
\begin{cases}
\dot x_{i,\eta}(t)=\xi_{i,\eta}(t),\ &x_{i,\eta}(0)=x_{i}^{0}\\[1mm]
\dot \xi_{i,\eta}(t)=-B(t,x_{i,\eta}(t))\xi_{i,\eta}^{\perp}(t)-\frac{1}{N\varepsilon}\underset{j:j\neq i}{\sum}\nabla_{x} V_{\eta}(x_{i,\eta}(t)-x_{j,\eta}(t)), \,\,\, &\xi_{i,\eta}(0)=\xi_{i}^{0}. 
\end{cases}
\label{regularized system}
\end{align}
Our main objective is to obtain a lower bound on the separation of the particles, {\it i.e.,} the quantity $\underset{i\neq j}{\min}\left\vert x_{i,\eta}(t)-x_{j,\eta}(t)\right\vert$. 
Using Theorem \ref{periodiccauchy lip thm} in the appendix, we can show that there exist a  global solution $(x_{1,\eta},\cdots\!,x_{N,\eta},\xi_{1,\eta},\cdots\!,\xi_{N,\eta})\in C^{1}(\mathbb{R}_{+};\mathbb{T}^{2N}\times \mathbb{R}^{2N})$ to \eqref{regularized system}. Note that 
\begin{align*}
\nabla_{x} V_{\eta}(x)=\nabla_{x}V_{\mathrm{loc}}(x)-2\Psi'(\frac{\left\vert x\right\vert^{2}}{\eta^{2}})\frac{x}{\eta^{2}}V_{\mathbb{R}^{2}}(x)+(1-\psi_{\eta}(x))\nabla_{x} V_{\mathbb{R}^{2}}(x). \end{align*}
Moreover, observe the estimates: 
\begin{align}
\left\vert (1-\psi_{\eta}(x))\nabla_{x} V_{\mathbb{R}^{2}}(x)\right\vert\lesssim \frac{1+\eta}{\eta}, \label{estimate for cutoff}       
\end{align}
and 
\begin{align}
\Big\vert 2\Psi'(\frac{\left\vert x\right\vert^{2}}{\eta^{2}})\frac{x}{\eta^{2}}V_{\mathbb{R}^{2}}(x)\Big\vert
\lesssim -\frac{\log(\eta)}{\eta}\lesssim \frac{1}{\eta^{2}}. \label{est for cutoff 2}      
\end{align}
We proceed by estimating $\left|x_{i,\eta}(t)-x_{i}^{0}\right|$. Integrating in time the equation for $x_{i}$ in \eqref{regularized system} and using Remark \ref{rem about conse of energy},
we obtain 
\begin{align}
\left\vert x_{i,\eta}(t)-x_{i}^{0}\right\vert 
&\leq \int_{0}^{t}\left\vert \xi_{i,\eta}(\tau)\right\vert\ \dd \tau  
\leq \Big(\int_{0}^{t}\left\vert \xi_{i}(\tau)\right\vert^{2}\ \dd\tau\Big)^{\frac{1}{2}}
\Big(\int_{0}^{t} \dd \tau\Big)^{\frac{1}{2}}\notag\\
&\leq \sqrt{T}\Big(2N\int_{0}^{t}\mathcal{F}_{\varepsilon,\eta,N}(\tau)\ \dd \tau\Big)^{\frac{1}{2}}
\leq \sqrt{2N\mathcal{F}_{\varepsilon,N}(0)}\,T, \label{6.4a}    
\end{align}
where we have denoted by $\mathcal{F}_{\varepsilon,\eta,N}(t)$ 
the total energy of the regularized system \eqref{regularized system}, defined with the non-negative kernel $V_\eta+C_0$,
and have used the fact that it is bounded by the initial energy of the original 
system \eqref{trajectories intro} because $V_{\eta}\leq V$. 
Thanks to \eqref{6.4a},
we have the bound 
\begin{align}
\left\vert x_{i,\eta}(t)-x_{j,\eta}(t)\right\vert&\geq \left\vert x_{i}^{0}-x_{j}^{0}\right\vert-\left\vert x_{i}^{0}-x_{i,\eta}(t)\right\vert-\left\vert x_{j}^{0}-x_{j,\eta}(t)\right\vert \notag\\
&\geq \underset{i\neq j}{\min}\left\vert x_{i}^{0}-x_{j}^{0}\right\vert-2\sqrt{2N\mathcal{F}_{\varepsilon,N}(0)}T.  \label{bound on distance from below}\end{align}
Now,  choose $T_{\ast}>0$ and $\eta>0$ as follows:  
\begin{itemize}
\item[{\rm(a)}]$\eta=\frac{1}{2}\underset{i\neq j}{\min}\left\vert x_{i}^{0}-x_{j}^{0}\right\vert;$
\item[{\rm(b)}] $T_{\ast}\leq \frac{\min_{i\neq j}\left\vert x_{i}^{0}-x_{j}^{0}\right\vert }{4\sqrt{2N\max\{1,\mathcal{F}_{\varepsilon,N}(0)\}}}.$  
\end{itemize}
In view of \eqref{bound on distance from below}, it follows that 
\begin{align*}
\min_{i\neq j}\left\vert x_{i,\eta}(t)-x_{j,\eta}(t)\right\vert\geq\eta\qquad \mbox{for all}\ t\in [0,T_{\ast}].      
\end{align*}
Since $\nabla_{x}V_{\eta}(x)\equiv\nabla_{x}V(x)$ on $\left\vert x\right\vert \geq \frac{\eta}{2}$, 
it follows that, for all $t\in [0,T_{\ast}]$,
\begin{align*}
\nabla_{x}V_{\eta}(x_{i,\eta}(t)-x_{j,\eta}(t))=\nabla_{x}V(x_{i,\eta}(t)-x_{j,\eta}(t))    
\end{align*}
which shows that 
$(x_{1,\eta}(t),\cdots\!,x_{N,\eta}(t),\xi_{1,\eta}(t),\cdots\!,\xi_{N,\eta}(t))$ is a solution to \eqref{trajectories intro} on $[0,T_{\ast}]$. On the open set where all pairwise distances exceed $\frac{\eta}{2}$, the vector field is locally Lipschitz; hence this solution is unique on $[0,T_*]$.
\end{proof}

\begin{lem}\label{lemma:6.2}
Let $X_{N}^{0}\in  \Delta_{N}^{c}$ and $\Xi_{N}^{0}\in \mathbb{R}^{2N}$, 
and let $(X_{N}(t),\Xi_{N}(t))\in C^{1}([0,T];\mathbb{T}^{2N}\times \mathbb{R}^{2N})$ 
be a solution to \eqref{trajectories intro} on $[0,T]$. Then
\begin{enumerate}
\item[{\rm(i)}] There exists  some constant  
$\delta=\delta(\varepsilon,N,\mathcal{F}_{\varepsilon,N}(0))>0$ such that 
\begin{align}
\min_{i\neq j}\left\vert x_{i}(t)-x_{j}(t)\right\vert\geq \delta 
\qquad\mbox{\ for all} \ t\in [0,T].  \label{bound from below on separation}    
\end{align}
\item[{\rm(ii)}] There exists 
$H_{0}=H_{0}(\varepsilon,T,\left\vert \Xi_{N}^{0}\right\vert,\delta )$ such that
\begin{align}
\underset{1\leq i\leq N}{\sup}\underset{t\in [0,T]}{\sup}\left\vert \xi_{i}(t)\right\vert\leq H_{0}. \label{H0 est} 
\end{align}
\end{enumerate}
\end{lem}

\begin{align}
\underset{1\leq i\leq N}{\sup}\underset{t\in [0,T]}{\sup}\left\vert \xi_{i}(t)\right\vert\leq H_{0}.  
\end{align}

\begin{thm}\label{periodiccauchy lip thm}
Suppose that $(X_{N}^{0},\Xi_{N}^{0})\in \mathbb{T}^{2N}\times \mathbb{R}^{2N}$, 
$V\in W^{2,\infty}(\mathbb{R}^{2})$, and $B\in C(\mathbb{R}_{+};W^{1,\infty}(\mathbb{R}^{2})),$ 
where $V$ and $B(t,\cdot)$ are $\mathbb{Z}^{2}$-periodic.
\begin{enumerate}
\item[{\rm(i)}] There exist a unique solution $(X_{N}(t),\Xi_{N}(t))\in C^{1}(\mathbb{R}_{+};\mathbb{R}^{2N}\times \mathbb{R}^{2N})$ of the Cauchy problem{\rm :} 
\begin{align}
\begin{cases}
\dot x_{i}(t)=\xi_{i}(t),\\
\dot \xi_{i}(t)=-B(t,x_{i}(t))\xi_{i}^{\perp}(t)-\frac{1}{N\varepsilon}\underset{j:j\neq i}{\sum}\nabla_{x} V(x_{i}(t)-x_{j}(t)),\\
x_{i}(0)=x_{i}^{0}, \quad \xi_{i}(0)=\xi_{i}^{0}. 
\end{cases}
\label{ode system appendix}
\end{align}
\item[{\rm(ii)}] If $(X_{N}(t),\Xi_{N}(t))$ is the unique global solution to \eqref{ode system appendix}, define $(\widetilde{X}_{N}(t),\widetilde{\Xi}_{N}(t)):= (X_{N}(t)\mod 1,\Xi_{N}(t))$. Moreover, let $(\widetilde{B},\widetilde{V})$ be the restriction of $(B,V)$ 
to $\mathbb{T}^{2}$, respectively.  
Then $(\widetilde{X}_{N}(t),\widetilde{\Xi}_{N}(t))\in C^{1}(\mathbb{R}_{+};\mathbb{T}^{2N}\times \mathbb{R}^{2N})$  is the unique solution to the Cauchy problem{\rm:} 
\begin{align}
\begin{cases}
\dot x_{i}(t)=\xi_{i}(t),\\
\dot \xi_{i}(t)=-\widetilde{B}(t,x_{i}(t))\xi_{i}^{\perp}(t)-\frac{1}{N\varepsilon}\underset{j:j\neq i}{\sum}\nabla_{x} \widetilde{V}(x_{i}(t)-x_{j}(t)), \\
\ x_{i}(0)=x_{i}^{0},\quad \xi_{i}(0)=\xi_{i}^{0}. 
\end{cases}
\label{ode system appendix torus}
\end{align}
\end{enumerate}
\end{thm}

\begin{thm}
Suppose that $X_{N}^{0}=(x_{1}^{N,0},\cdots\!,x_{N}^{N,0})\in \Delta_{N}^{c}$,
$\Xi_{N}^{0}=(\xi_{1}^{N,0},\cdots\!,\xi_{N}^{N,0})\in \mathbb{R}^{2N}$, 
and $B\in C(\mathbb{R}_{+};W^{1,\infty}(\mathbb{T}^{2}))$. 
Then there exists a unique global solution $(X_{N}(t),\Xi_{N}(t))\in C^{1}(\mathbb{R}_{+}; \mathbb{T}^
{2N}\times \mathbb{R}^{2N})$ to \eqref{trajectories intro}.    
\end{thm}

\begin{proof}  We divide into two steps.

\smallskip
\textbf{1. Existence}.
By 
Theorem \ref{Short time well posed}, there exist some $T_{\ast}>0$ and a solution $(X_{N},\Xi_{N})\in C^{1}([0,T_{\ast}];\mathbb{T}^{2N}\times \mathbb{R}^{2N})$ to \eqref{trajectories intro}. We proceed by induction. Suppose that we have a solution $(X_{N},\Xi_{N})\in C^{1}([0,nT_{\ast}];\mathbb{T}^{2N}\times \mathbb{R}^{2N})$ for \eqref{trajectories intro}. Consider the  Cauchy problem 
\begin{align}
\begin{cases}
\dot y_{i}(t)=\lambda_{i}(t),\\[1mm]
\dot \lambda_{i}(t)=-B(t+nT_{\ast},y_{i}(t))\lambda_{i}^{\perp}(t)-\frac{1}{N\varepsilon}\underset{j:j\neq i}{\sum}\nabla_{x}V(y_{i}(t)-y_{j}(t)),\\
\ y_{i}(0)=x_{i}(nT_{\ast}),\quad \lambda_{i}(0)=\xi_{i}(nT_{\ast}). 
\end{cases} \label{iterative step}   
\end{align}
By Lemma \ref{lemma:6.2}, we have
\begin{align*}
\min_{i\neq j}\left\vert y_{i}(0)-y_{j}(0)\right\vert=\min_{i\neq j}\left\vert x_{i}(nT_{\ast})-x_{j}(nT_{\ast})\right\vert\geq \delta(\varepsilon,N,\mathcal{F}_{\varepsilon,N}(0)).      
\end{align*}
Therefore, by Proposition \ref{Short time well posed}, 
it follows that there exists a solution 
$(Y_{N},\Lambda_{N})\in C^{1}([0,T_{\ast}];\mathbb{T}^{2N}\times \mathbb{R}^{2N})$ 
to \eqref{iterative step} 
(note that we can take the same $T_{\ast}$ because it depends only on the initial separation $\delta$ and the initial energy $\mathcal{F}_{\varepsilon,N}(0)$). Set 
\begin{align*}
Z_{N}(t):=\begin{cases}
X_{N}(t) &\quad \mbox{for $t\in [0,nT_{\ast}]$},\\
Y_{N}(t-nT_{\ast})&\quad \mbox{for $t\in [nT_{\ast},(n+1)T_{\ast}]$},
\end{cases}
\end{align*}
and 
\begin{align*}
\Pi_{N}(t):=\begin{cases}
\Xi_{N}(t)&\quad\mbox{for $t\in [0,nT_{\ast}]$},\\
\Lambda_{N}(t-nT_{\ast}) &\quad\mbox{for $t\in [nT_{\ast},(n+1)T_{\ast}]$}. 
\end{cases}    
\end{align*}
Then it is readily checked that 
\begin{itemize}
\item[{\rm(a)}] $(Z_{N},\Pi_{N})\in C^{1}([0,(n+1)T_{\ast}];\mathbb{T}^{2N}\times \mathbb{R}^{2N});$
\smallskip
\item[{\rm(b)}] $(Z_{N},\Pi_{N})$ is a solution to \eqref{trajectories intro} on $[0,(n+1)T_{\ast}]$.
\end{itemize}
This concludes the induction step, and hence it follows that, for every $n\in \mathbb{N}$, 
there exists a solution on $[0,nT_{\ast}]$ to \eqref{trajectories intro}. 

\smallskip
\textbf{2. Uniqueness}. Let $(X_{N},\Xi_{N})\in C^{1}([0,T];\mathbb{T}^{2N}\times \mathbb{R}^{2N})$ and $(Y_{N},\Lambda_{N})\in C^{1}([0,T];\mathbb{T}^{2N}\times \mathbb{R}^{2N})$ be two solutions of \eqref{trajectories intro} with the same initial data $(X_{N}^{0},\Xi_{N}^{0})$. Denote 
\begin{align*}
d_{\min}=\min\Big\{\underset{t\in [0,T]}{\min}\underset{i\neq j}{\min}\left\vert x_{i}(t)-x_{j}(t)\right\vert,\underset{t\in[0,T]}{\min}\min_{i\neq j} \left\vert y_{i}(t)-y_{j}(t)\right\vert\Big\}.     
\end{align*}
Note that $(X_{N},\Xi_{N})$ and $(Y_{N},\Lambda_{N})$ satisfy the truncated equation \eqref{regularized system} with $\eta=d_{\min}$. 
We write system \eqref{regularized system} concisely as 
\begin{align*}
\frac{\dd}{\dd t}(X_{N}(t),\,\Xi_{N}(t))
=(F_{1}(t,X_{N}(t),\Xi_{N}(t)),\,F_{2}(t,X_{N}(t),\Xi_{N}(t)))  \end{align*}
where $F_{1}:\mathbb{R}_{+}\times \mathbb{T}^{2N}\times \mathbb{R}^{2N}\rightarrow \mathbb{R}^{2N}$ and $F_{2}:\mathbb{R}_{+}\times \mathbb{T}^{2N}\times \mathbb{R}^{2N}\rightarrow \mathbb{R}
^{2N}$ are the vector fields defined by 
\begin{align*}
F_{1}(t,X_{N},\Xi_{N}):= (\xi_{1},\cdots\!,\xi_{N}),    
\end{align*}
and 
\begin{align*}
F_{2}(t,X_{N},\Xi_{N}):=\Big(-\xi_{i}^{\perp}B(t,x_{i})-\frac{1}{N\varepsilon}\sum_{j:j\neq i}\nabla_{x}V_{\eta}(x_{i}-x_{j})\Big)_{i=1}^{N}.
\end{align*}
We have 
\begin{align*}
&\frac{\dd}{\dd t}(\left\vert X_{N}(t)-Y_{N}(t)\right\vert^{2}+\left\vert \Xi_{N}(t)-\Lambda_{N}(t)\right\vert^{2})\\
&=2(X_{N}(t)-Y_{N}(t))\cdot (F_{1}(t,X_{N}(t),\Xi_{N}(t))-F_{1}(t,Y_{N}(t),\Lambda_{N}(t)))\\
&\quad+2(\Xi_{N}(t)-\Lambda_{N}(t))\cdot (F_{2}(t,X_{N}(t),\Xi_{N}(t))-F_{2}(t,Y_{N}(t),\Lambda_{N}(t))):= \mathcal{T}_{1}+\mathcal{T}_{2}.    \end{align*}
For brevity, we omit the time variable in what follows. We estimate $\mathcal{T}_{1}$:\begin{align}
\mathcal{T}_{1}\leq 2\left\vert X_{N}-Y_{N}\right\vert\left\vert \Xi_{N}-\Lambda_{N}\right\vert\leq \left\vert X_{N}-Y_{N}\right\vert^{2}+\left\vert \Xi_{N}-\Lambda_{N}\right\vert^{2}.  \label{T1 estimate}        
\end{align}
To estimate $\mathcal{T}_{2}$,  start by noting that 
\begin{align*}
|\mathcal{T}_{2}|&\leq 2\left\vert \Xi_{N}-\Lambda_{N}\right\vert \left\vert F_{2}(X_{N},\Xi_{N})-F_{2}(Y_{N},\Lambda_{N})\right\vert\\
&\leq 2\left\vert \Xi_{N}-\Lambda_{N}\right\vert \left\vert F_{2}(X_{N},\Xi_{N})-F_{2}(X_{N},\Lambda_{N})\right\vert+2\left\vert \Xi_{N}-\Lambda_{N}\right\vert \left\vert F_{2}(X_{N},\Lambda_{N})-F_{2}(Y_{N},\Lambda_{N})\right\vert\\
&:=\mathcal{T}_{2}^{1}+\mathcal{T}_{2}^{2}. 
\end{align*}
We have 
\begin{align*}
\left\vert F_{2}(X_{N},\Xi_{N})-F_{2}(X_{N},\Lambda_{N})\right\vert
\leq \left\Vert B\right\Vert_{\infty}\big\vert (\xi_{1}^{\perp},\cdots\!,\xi_{N}^{\perp})-(\lambda_{1}^{\perp},\cdots\!,\lambda_{N}^{\perp})\big\vert
=\left\Vert B\right\Vert_{\infty}\left\vert \Xi_{N}-\Lambda_{N}\right\vert,          
\end{align*}
so that
\begin{align}
\mathcal{T}_{2}^{1}\leq  2\left\Vert B\right\Vert_{\infty} \left\vert \Xi_{N}-\Lambda_{N}\right\vert^{2}.\label{est on T12}     
\end{align}
We introduce the notation 
\begin{itemize}
  \item[{\rm(a)}]$L=L(d_{\min})=\mathrm{Lip}(\nabla_{x}V_{\eta})$.
   \item[{\rm(b)}] $H_{0}$ as in \eqref{H0 est}. 
\end{itemize}
With this notation, we estimate 
\begin{align*}
&\left\vert F_{2}(X_{N},\Lambda_{N})-F_{2}(Y_{N},\Lambda_{N})\right\vert\\
&\leq H_{0}\left\Vert \nabla_xB\right\Vert_{L^{\infty}([0,T]\times\mathbb T^2)}\left\vert X_{N}-Y_{N}\right\vert
+\frac{L}{N\varepsilon}
\bigg(\sum_{i=1}^{N}\Big(\sum_{j\ne i}
\bigl(|x_i-y_i|+|x_j-y_j|\bigr)\Big)^2\bigg)^{1/2}\\
&\leq \big(H_{0}\left\Vert \nabla_xB\right\Vert_{L^{\infty}([0,T]\times\mathbb T^2)}+\frac{2L}{\varepsilon}\big)\left\vert X_{N}-Y_{N}\right\vert.
\end{align*}
It follows that 
\begin{align}
\mathcal{T}^{2}_{2}\leq \big(H_{0}\left\Vert B\right\Vert_{L^{\infty}_{t}W^{1,\infty}_{x}} +\frac{2L}{\varepsilon}\big)\big(\left\vert X_{N}-Y_{N}\right\vert^{2}+\left\vert \Xi_{N}-\Lambda_{N}\right\vert^{2}\big). \label{est on T22} 
\end{align}
Gathering \eqref{est on T12}--\eqref{est on T22}, we have
\begin{align}
\mathcal{T}_{2}
\leq \big(H_{0}\left\Vert B\right\Vert_{L^{\infty}_{t}W^{1,\infty}_{x}}+\frac{2L}{\varepsilon}+2\left\Vert B\right\Vert_{\infty}\big)\big(\left\vert X_{N}-Y_{N}\right\vert^{2}+\left\vert \Xi_{N}-\Lambda_{N}\right\vert^{2}\big). \label{T2 estimate}  
\end{align}
To conclude, combining \eqref{T1 estimate}--\eqref{T2 estimate}  yields the following estimate 
for some constant $C>0$: 
\begin{align*}
\frac{\dd}{\dd t}(\left\vert X_{N}(t)-Y_{N}(t)\right\vert^{2}+\left\vert \Xi_{N}(t)-\Lambda_{N}(t)\right\vert^{2})\leq C(\left\vert X_{N}(t)-Y_{N}(t)\right\vert^{2} +\left\vert \Xi_{N}(t)-\Lambda_{N}(t)\right\vert^{2} ).     
\end{align*}
Then the uniqueness follows directly by Gr\"onwall's lemma. 
\end{proof}

\end{document}
